% III. Експериментална проверка и резултати. Слети раздели V и VI.
% Всяко число тук е от изпълнение на quorum bench или quorum demo. Нищо не е
% пренесено по памет.
\section{Експериментална проверка и резултати}
\label{sec:results}

\subsection{Методика и среда}

Проверката отговаря на четири въпроса. Дава ли системата правилен отговор там,
където той е известен предварително. Съвпадат ли класиранията на двата метода за
оценяване и достига ли разминаването до портфейла. Каква е цената на точното
решаване спрямо евристиката. Колко устойчива е препоръката спрямо изменение на
допусканията. Използвани са две групи сценарии. Първата е породена от генератор с фиксирано
начално състояние, с контролиран брой проекти, две измерения на капацитета,
предхождане с вероятност една четвърт и няколко двойки взаимно изключване;
бюджетът е на четиридесет процента от исканото, за да е ограничението действащо.
Втората е един съставен сценарий, \code{scenarios/example.json}, с дванадесет
проекта, четири критерия, два ресурса и бюджет, покриващ около половината от
исканото. Трета група, сценарий по действителна практика в организация, беше
предвидена, но такъв не беше получен и групата отпада. Тя не се замества с
измислен пример: сценарий, съчинен така, че да изглежда действителен, не носи
повече доказателствена сила от породения по правило, а внушава, че носи.

Всички числа са получени на една и съща машина: AMD Ryzen 5 3600, 16 GiB,
Windows 11 Pro N, g++ 15.2.0 (MinGW-w64), \code{-std=c++20}, режим \code{Release},
CMake 4.3.2 и Ninja 1.13.2, без външни зависимости. Времето се измерва вътре в
програмата около самия алгоритъм, при едно и също нишково изпълнение.
Измерванията се изпълняват с \code{quorum bench}, а рискът и чувствителността с
\code{quorum demo}; и двете не приемат параметри, които биха променили отчетените
числа. Условие за валидност: точното решение не може да бъде по-лошо от
евристичното, а нарушение спира експеримента, вместо да се отчете като резултат.

\subsection{Препоръката и защо всеки проект е вътре или вън}

При бюджет 1150 и искане 2106 системата финансира седем от дванадесетте проекта
за 1123, тоест 97,7 процента от бюджета. Инженерният ресурс е изчерпан напълно,
26,0 от 26,0, докато операционният остава на 9,0 от 11,0. Търсенето обхожда 357
върха и доказва оптималност за 0,031 ms; лакомата евристика достига 97,55 процента
от оптималната стойност за 0,006 ms.

Критичният път минава през P03 и P04 и дава обща продължителност 34,7 периода.
Останалите пет финансирани проекта имат запас между 16,8 и 22,8 периода, тоест
забавяне в който и да е от тях не измества края на програмата.

Отхвърлянията идват с причина. И петте отхвърлени проекта се спират от бюджета:
остатъкът е 27, а цените им са между 62 и 325. Сред тях е P02, класиран първи и
по двата метода за оценяване. Предшественикът му P01 е финансиран, взаимно
изключване не го засяга, спира го единствено цената. Това е случаят, в който
подредбата по оценка и съставът на портфейла се разминават и системата дължи на
читателя причината, а не само списъка.

\subsection{Цената на решение, взето върху очаквани стойности}

При двадесет хиляди извадки върху избрания портфейл общата стойност има средна
1162 и стандартно отклонение 38,6, с деветдесетпроцентов централен интервал
$[1100; 1227]$. Срокът има средна 37,35 периода при стандартно отклонение 3,32 и
интервал $[32{,}02; 43{,}04]$. Оттук следва основният резултат на проекта:
портфейлът се побира в бюджета по очаквани стойности, 1123 срещу 1150, но
вероятността действителната сума да остане в бюджета е \textbf{38,7 процента},
тоест около три на пет шанс за преразход. Това е числото, което таблицата с
очаквани стойности не показва, и е причината анализът на риска да е част от
системата, а не допълнение към нея.

Разминаването между 1123 и 1162 е 3,5 процента и е следствие от съзнателно
избраното разпределение, а не от грешка: подборът работи с очаквани стойности по
PERT, при които модата получава тегло четири, а извадката се тегли от триъгълно
разпределение със средна, равна на средно аритметично на трите точки, а оценките
в сценария са несиметрични надясно. Средният срок надхвърля детерминираните 34,7
по същата причина, плюс ефекта на сливането на пътищата.

Броят извадки по подразбиране не е избран по усещане: оценката за p95 е повторена
с десет начални състояния на генератора и размахът пада от 7,99 при хиляда
извадки до 1,20 при сто хиляди.

\subsection{Устойчивост на препоръката}

При изменение на теглата, шестнадесет изпълнения, пет позиции остават в портфейла
винаги: P03, P04, P07, P09 и P11. P01 и P08 излизат по веднъж, P06 и P10 влизат
само по веднъж, а P02, P05 и P12 не влизат никога. При изменение на бюджета от
шестдесет до сто и четиридесет процента, девет изпълнения, картината е
по-подвижна и е дадена на Фиг.~\ref{fig:stability}.

\begin{figure}[H]
\centering
\begin{tikzpicture}[x=1.08cm, y=0.5cm]
  \foreach \r/\name in {1/P01,2/P02,3/P03,4/P04,5/P05,6/P06,7/P07,8/P08,9/P09,10/P10,11/P11,12/P12}
    \node[anchor=east, font=\small] at ({0.45}, {-\r}) {\name};
  \foreach \c/\f in {1/{0,6},2/{0,7},3/{0,8},4/{0,9},5/{1,0},6/{1,1},7/{1,2},8/{1,3},9/{1,4}}
    \node[anchor=south, font=\small] at ({\c}, {-0.6}) {\f};
  \foreach \r in {1,...,12}
    \foreach \c in {1,...,9}
      \draw[gray!35, line width=0.3pt] ({\c-0.42}, {-\r-0.2}) rectangle ({\c+0.42}, {-\r+0.2});
  \foreach \c/\r in {1/1,2/1,3/1,5/1,6/1,7/1,8/1,9/1,
                     2/2,3/2,
                     4/3,5/3,6/3,7/3,8/3,9/3,
                     4/4,5/4,6/4,7/4,8/4,9/4,
                     6/6,7/6,8/6,9/6,
                     1/7,2/7,3/7,4/7,5/7,6/7,7/7,8/7,9/7,
                     1/8,2/8,3/8,4/8,5/8,6/8,7/8,8/8,9/8,
                     1/9,2/9,3/9,4/9,5/9,6/9,7/9,8/9,9/9,
                     1/10,4/10,
                     3/11,4/11,5/11}
    \fill[blue!45] ({\c-0.42}, {-\r-0.2}) rectangle ({\c+0.42}, {-\r+0.2});
  \node[anchor=north, font=\small] at ({5}, {-12.8}) {бюджет като дял от зададения};
\end{tikzpicture}
\caption{Позиции в избрания портфейл при изменение на бюджета. Запълнена клетка означава финансиран проект.}
\label{fig:stability}
\end{figure}

Три позиции, P07, P08 и P09, са в портфейла при всеки от деветте бюджета. Две,
P05 и P12, не влизат при нито един, включително при сто и четиридесет процента:
те губят по оценка, а не само по бюджет. Между тях са седемте подвижни позиции.

Фигурата показва и нещо, което делът в проценти би скрил: присъствието на проект
не е монотонно по бюджета. P01 е финансиран при шестдесет, седемдесет и осемдесет
процента, отпада при деветдесет и се връща при сто; P10 е финансиран при
шестдесет и при деветдесет, но не между тях. Причината е, че при по-широк бюджет
става възможна двойката P03 и P04, която измества други позиции: сумарната
стойност расте монотонно, но съставът не може. Оттук следва предупреждението за
управленския читател. Изречение като „P01 е препоръчан в осемдесет и девет
процента от случаите“ е вярно, но подвеждащо, защото внушава интервал, който не
съществува. Прочитът изисква решетката, а не само процента, затова системата
отпечатва и двете.

\subsection{Двата метода за оценяване дават един и същ портфейл}

Теглата, изведени от матрицата на попарните сравнения, са $0{,}1891$ за
стратегическото съответствие, $0{,}3509$ за възвръщаемостта, $0{,}3509$ за
регулаторната необходимост и $0{,}1091$ за риска, при коефициент на съгласуваност
$0{,}0036$, далеч под прага $0{,}10$. При тези тегла двата метода класират
дванадесетте проекта различно на шест позиции: P09, P06 и P04 излизат по-напред
по TOPSIS съответно с една, две и една позиция, а P10, P12 и P03 по-назад с две,
една и една. Коефициентът на рангова корелация на Spearman е $0{,}9580$; всички
разминавания са между пето и десето място.

Отражението върху решението е нулево: изпълнен поотделно с двата вектора от
оценки, подборът връща един и същ портфейл, P01, P03, P04, P07, P08, P09 и P11.
При този сценарий изборът между двата метода не е решаващ, и това е по-полезно за
управленския читател от която и да е от двете подредби поотделно. Обратното, че
методът никога не е от значение, не следва оттук: проверено е върху един сценарий.

\subsection{Коректност и цена на точното решаване}

Стойността от разклоняването и ограничаването е сравнена с изчерпателно изброяване
на всички подмножества при 8, 12, 16 и 18 проекта, по двадесет задачи на размер.
Съвпадението е пълно при всичките осемдесет. Средното време на задача расте от
0,0039 до 0,1069 ms, докато изброяването расте от 0,021 до 32,353 ms; при двадесет
и пет проекта еталонът би отнел часове, което е и горната граница на проверката.

\begin{table}[H]
\centering
\caption{Разклоняване и ограничаване срещу лакома евристика, по пет задачи на размер}
\label{tab:exact-greedy}
\small
\begin{tabular}{@{}cccccc@{}}
\toprule
\textbf{Проекти} & \textbf{Точно} & \textbf{Лакомо} & \textbf{Дял, \%} & \textbf{Върхове} & \textbf{Точно, ms} \\
\midrule
10 & 0,5753 & 0,5197 & 89,87 & 74     & 0,008 \\
14 & 0,7875 & 0,6940 & 87,84 & 225    & 0,018 \\
18 & 1,1170 & 0,9141 & 81,11 & 920    & 0,077 \\
22 & 1,2272 & 1,0708 & 87,43 & 3440   & 0,237 \\
26 & 1,5857 & 1,2652 & 79,01 & 7937   & 0,753 \\
30 & 1,8450 & 1,6092 & 86,58 & 18439  & 1,710 \\
34 & 2,2365 & 1,8923 & 84,75 & 54127  & 5,530 \\
38 & 2,4237 & 2,1413 & 88,53 & 73295  & 7,862 \\
\bottomrule
\end{tabular}
\end{table}

Евристиката губи между десет и двадесет и един процента от оптималната стойност,
без тенденция с размера, което е характерно за лаком избор при действащи
ограничения за предхождане. Цената на точността е малка: при тридесет и осем
проекта търсенето отнема под осем милисекунди, тоест анализът на чувствителността
остава под секунда и при задачи, много по-големи от съставния сценарий.
Отношението между двата алгоритъма обаче расте от около единица при десет проекта
до около триста и петдесет при тридесет и осем, тоест границата на точното
решаване е в порядъка на десетки проекти.
